% META: {"id": "TCOMPL-MF-EX-011", "chapitre": "TCOMPL-MODELES-FONCTION", "type_objet": "exercice", "capacites": ["TCOMPL-MF-C6"], "capacites_codes": ["C6"], "methodes": ["M6"], "parcours": 2, "competences": ["representer", "calculer", "raisonner"], "duree_min": 22, "mode_creation": "ex_nihilo", "sources_inspiration": [], "fichier_tex": "chapitres/TCOMPL-MODELES-FONCTION/exercices/TCOMPL-MF-EX-011.tex", "corrige_tex": "chapitres/TCOMPL-MODELES-FONCTION/corriges/TCOMPL-MF-CO-011.tex", "authoring": {"PEDAGOGICAL_GAP_ID": "GAP-3C9D7BB4136C", "CAPACITY_ID": "C6", "EXERCISE_ROLE": "CONTEXT_VARIATION", "DIFFICULTY": 2, "AUTHORING_REASON": "le point moyen n'etait calcule que pour lui-meme, sur une liste de couples hors contexte ; rien ne montrait a quoi il sert, alors que c'est le point par lequel tout ajustement affine doit passer", "ORIGINALITY_PROVENANCE": "ex_nihilo"}, "status": "generated"}
% BEGIN-VERIFY
% from sympy import Rational, Point2D, symbols, Eq, solve, simplify
% xs = [2, 5, 8, 11, 14, 20]
% ys = [310, 266, 240, 196, 158, 90]
% assert len(xs) == len(ys) == 6
% assert sum(xs) == 60 and sum(ys) == 1260
% xbar, ybar = Rational(sum(xs), 6), Rational(sum(ys), 6)
% assert xbar == 10 and ybar == 210
% g1 = Point2D(Rational(sum(xs[:3]), 3), Rational(sum(ys[:3]), 3))
% g2 = Point2D(Rational(sum(xs[3:]), 3), Rational(sum(ys[3:]), 3))
% assert g1 == Point2D(5, 272) and g2 == Point2D(15, 148)
% assert Point2D(xbar, ybar) == g1.midpoint(g2)
% pente = (g2.y - g1.y) / (g2.x - g1.x)
% assert pente == Rational(-62, 5)
% x = symbols('x')
% droite = pente * (x - xbar) + ybar
% assert simplify(droite - (Rational(-62, 5) * x + 334)) == 0
% assert droite.subs(x, 10) == 210
% assert droite.subs(x, 17) == Rational(616, 5)
% assert float(droite.subs(x, 17)) == 123.2
% assert droite.subs(x, g1.x) == g1.y and droite.subs(x, g2.x) == g2.y
% END-VERIFY
\begin{exercice}{TCOMPL-MF-EX-011}{2}{22}
Un gestionnaire d'immeuble relève, pour six mois de l'année, la température
moyenne extérieure $x_i$ (en degrés Celsius) et la consommation de chauffage
$y_i$ (en kilowattheures par logement) :
\[
  (2,310),\ (5,266),\ (8,240),\ (11,196),\ (14,158),\ (20,90).
\]
\begin{enumerate}
  \item Représenter le nuage de points dans un repère, en prenant $1$~cm pour
    $2\,^\circ$C en abscisse et $1$~cm pour $40$~kWh en ordonnée.
  \item Calculer les coordonnées du point moyen $G$ et le placer sur le graphique.
  \item On partage le nuage en deux groupes de trois points, dans l'ordre des
    abscisses croissantes. Calculer les points moyens $G_1$ et $G_2$ de ces deux
    groupes. Vérifier que $G$ est le milieu de $[G_1G_2]$, et expliquer pourquoi
    ce n'est pas un hasard ici.
  \item Déterminer une équation de la droite $(G_1G_2)$, appelée droite de Mayer.
    Vérifier qu'elle passe par $G$.
  \item Estimer la consommation pour une température moyenne de $17\,^\circ$C.
  \item Le gestionnaire propose une autre droite d'ajustement, d'équation
    $y=-13x+340$. Passe-t-elle par $G$ ? Que peut-on en dire ?
\end{enumerate}
\end{exercice}
